\section{What a Thing Is, and That It Is}

The world that inherited the Greek instruments professed, on entirely non-philosophical grounds, that the whole of things had been called out of nothing. Whether that profession is true is not the present question, and nothing in this section will assume it. But it explains a change of focus. Thinkers who confessed creation had a standing reason to ask the question Athens never isolated---what accounts for there being anything to move?---and to ask it with Greek rigor. The thirteenth century, and Thomas Aquinas in particular, brought the two inheritances into one argument.

\subsection{A distinction quiet enough to miss}

Aquinas's decisive move is a distinction so quiet it can be missed. Consider any definable thing---a horse, a hexagon, a helium atom. Everything that belongs to \emph{what it is} can be laid out in the definition, and the definition will be complete without ever mentioning \emph{that it is}. You can know exhaustively what a thing would be while its existence remains an open question; naturalists have written precise descriptions of animals that turned out not to exist, and physicists of particles that turned out to exist after all. In no case did learning the answer add a new clause to the definition. Therefore what a thing is and that it is are really distinct in every such thing: essence does not include existence, any more than the blueprint of the house includes the house's standing there. A thing's nature is related to its existing as a capacity to what fills it---which is to say, in Aristotle's currency, as potency to act.

The distinction can be stated so briskly that it sounds like a verbal point, so it is worth watching what it is not. It is not the claim that existence is one more feature a thing might have or lack, alongside mass and color---a claim that would be false and that the tradition denies. Existence is not an item in the inventory; it is the condition of there being an inventory. Nor is it a claim about the order of discovery, as though we first meet natures and then check them against reality. It is a claim about the structure of any finite thing: the nature does not account for the being, and the being is not contained in the nature.

\subsection{\latin{De ente et essentia}: the argument at its source}

The fullest early statement is a short treatise Aquinas wrote as a young man, \emph{On Being and Essence}, and it is worth following in its own steps, because the compressed versions in the \emph{Summa} presuppose it. The argument begins from what the mind can do: ``every essence or quiddity can be understood without anything being understood about its existing; for I can understand what a man is, or what a phoenix is, and still not know whether it has being in the nature of things.''\endnote{Thomas Aquinas, \emph{De ente et essentia}, ch.~4 in the Leonine chapter division, corresponding to ch.~3 of the Baur--Koch text presented in the Corpus Thomisticum web corpus, where the passage was read: \latin{Omnis autem essentia vel quiditas potest intelligi sine hoc quod aliquid intelligatur de esse suo; possum enim intelligere quid est homo vel Phoenix et tamen ignorare an esse habeat in rerum natura.} Latin checked 2026-07-25 at corpusthomisticum.org (\texttt{oee.html}); the English renderings of this treatise in the present section are labeled working glosses governed by the quoted Latin, no public-domain English translation having been identified. The chapter-numbering discrepancy is real and is recorded here rather than silently resolved: editions divide the treatise differently, so citations to ``ch.~4'' and ``ch.~3'' of it frequently designate the same text.} The pairing of examples is exact. A man is something we know to exist; a phoenix, in Aquinas's day as in ours, is something about which the question stands open or has been closed in the negative. The definition is the same document in every case. Whatever the answer about existence turns out to be, nothing is added to or struck from the account of what the thing is---and that invariance is the whole evidence.

From there Aquinas moves in three steps. First: since existence is not included in what a thing is, existence must be \emph{added} to the nature---``therefore it is clear that existence is other than essence or quiddity, unless perhaps there is some thing whose quiddity is its own existence,'' and such a thing, he argues immediately, ``can only be one, and first.''\endnote{\latin{Ergo patet quod esse est aliud ab essentia vel quiditate, nisi forte sit aliqua res, cuius quiditas sit ipsum suum esse; et haec res non potest esse nisi una et prima}; same locus, checked 2026-07-25. The uniqueness argument turns on there being no way to multiply a thing that is nothing but existence: multiplication requires an added difference, a receiving matter, or a distinction between what is absolute and what is received, and \latin{esse tantum} admits none of these.} Second: whatever belongs to a thing without being part of its nature comes either from the thing's own principles or from outside---and here Aquinas reaches, in the very sentence that carries the argument, for an illustration this essay has been using since its first page: ``as risibility in man, or as the light in the air from the influence of the sun.''\endnote{\latin{Omne autem quod convenit alicui vel est causatum ex principiis naturae suae, sicut risibile in homine, vel advenit ab aliquo principio extrinseco, sicut lumen in aere ex influentia solis}; same locus, checked 2026-07-25. That the lit air appears here, at the joint where existence is shown to be received, and again at \emph{Summa theologiae} I, q.~8, a.~1 and I, q.~104, a.~1, is the reason this essay treats the image as the tradition's rather than as an illustration of its own devising.} Third, the exclusion: existence cannot be caused by the thing's own nature, ``for then some thing would be its own cause and would produce itself in being, which is impossible.''\endnote{\latin{Non autem potest esse quod ipsum esse sit causatum ab ipsa forma vel quiditate rei (dico sicut a causa efficiente) quia sic aliqua res esset sui ipsius causa et aliqua res seipsam in esse produceret, quod est impossibile}; same locus, checked 2026-07-25. Aquinas's parenthesis (``I mean as by an efficient cause'') is precise: he is not denying that a nature explains its own properties, only that it can confer its own existing.}

The conclusion follows in one move, and it is the move the rest of this essay depends on: ``therefore every such thing, whose existence is other than its nature, must have existence from another. And since everything that is through another is traced back to what is through itself as to a first cause, there must be some thing which is the cause of existing for all things, in that it is existence alone.''\endnote{\latin{Ergo oportet quod omnis talis res, cuius esse est aliud quam natura sua habeat esse ab alio. Et quia omne, quod est per aliud, reducitur ad illud quod est per se sicut ad causam primam, oportet quod sit aliqua res, quae sit causa essendi omnibus rebus, eo quod ipsa est esse tantum. Alias iretur in infinitum in causis, cum omnis res, quae non est esse tantum, habeat causam sui esse}; same locus, checked 2026-07-25. Note the two premises doing the work: that what is through another traces to what is through itself, and that an endless series of caused causes of \emph{being} explains nothing. The regress denied here is the same essentially ordered kind analyzed in section~4 below; the treatise makes no claim about the world's age.} A borrower of existence traces to an owner of existence, or the account never closes. And what the owner is, on this analysis, is not a very large thing that happens to exist without help. It is \latin{esse tantum}---existence, and nothing else besides.

\subsection{The schools that differ}

Honesty requires saying that the thirteenth century did not receive this argument as an obvious truth, and that the disagreement was internal to the tradition rather than an assault from outside it. Within a generation the question ``how are essence and existence distinguished in creatures?'' had become one of the most contested in the schools. Giles of Rome defended a real distinction in a stronger form, at times speaking of essence and existence as distinct \emph{things}---language that invited critics to picture existence as an entity rather than a principle. Henry of Ghent denied a real distinction while insisting on more than a merely conceptual one, proposing a third, ``intentional'' distinction between them. Godfrey of Fontaines rejected both and held essence and existence to be really identical, differing only in the way they signify.\endnote{The state of the thirteenth-century dispute is reported here at a stated secondary ceiling: John F. Wippel, ``Godfrey of Fontaines,'' \emph{Stanford Encyclopedia of Philosophy}, section~6 (``Essence and Existence''), substantive revision of 23 October 2018, read 2026-07-25 at plato.stanford.edu. That witness supplies the positions summarized: Giles's real distinction and its \latin{res} language, Henry's intentional distinction, and Godfrey's identity thesis (``for him they are identical, and differ only in the way they signify''). No text of Giles, Henry, or Godfrey was collated for this essay, and no claim is made about their arguments finer than the cited witness supports. The later scholastic and early-modern continuations of the dispute are not summarized here at all, and no position is attributed to schools whose texts were not checked.} The point of naming them is not to adjudicate. It is that a reader who finds the distinction less than self-evident is in respectable medieval company, and that this essay's adoption of the Thomistic position is an adoption, not a report of consensus.

A second and quite different objection belongs here too, because it is the one a contemporary reader is more likely to have met. Since Hume and Kant, and with new precision since Frege and Russell, a long line of philosophers has denied that existence is a property of individuals at all: to say that dinosaurs do not exist, on this account, is not to say something about certain individuals but to say that a certain concept is not instantiated.\endnote{Reported at a stated secondary ceiling: Michael Nelson, ``Existence,'' \emph{Stanford Encyclopedia of Philosophy}, substantive revision of 5 May 2020, read 2026-07-25 at plato.stanford.edu, which narrates both sides of this history---Aquinas's argument in \emph{De ente et essentia} ch.~4 (the entry's own citation) and the counter-tradition of Hume, Kant, Frege, and Russell for whom existence is not a first-order property of individuals. No primary text of Hume, Kant, Frege, or Russell was collated for this essay, and the paragraph above claims nothing finer than the cited witness supports. A full engagement with the quantificational account of existence is outside this essay's scope, and its absence is recorded in the appendix.} That tradition and this one are not simply colliding, and it is worth seeing why. The scholastic claim was never that existing is one of a thing's attributes---the previous subsection denied exactly that---but that a finite nature does not account for its own being in act. A reader persuaded by Frege about the logic of ``exists'' may still ask why this particular nature is in act rather than not, and that question, not a question of predicate logic, is the one the argument answers. The essay records the collision without pretending to settle it, and states plainly what it needs: that no finite nature accounts for its own existing.

\subsection{Participation, and the terminus}

Once the distinction is in hand, the Greek machinery engages it with a click. Nothing confers on itself what it does not have; a nature that does not include existence cannot be the source of its own existing. Existence must then be received---and received the way the song is received from the singer, not the way the house was received from the builder, because the receiving is not an event in the thing's past but the present condition of there being a thing. Aquinas says it in the vocabulary Plato minted: ``all beings apart from God are not their own being, but are beings by participation,'' and whatever has a perfection by participation is caused in it by what has that perfection essentially, ``as iron becomes ignited by fire.''\endnote{\emph{Summa theologiae} I, q.~44, a.~1, corpus (``Therefore all beings apart from God are not their own being, but are beings by participation''; \latin{omnia alia a Deo non sint suum esse, sed participant esse}); wording checked 2026-07-25, English at New Advent, Latin at Corpus Thomisticum. The iron-and-fire comparison is Aquinas's own in the same corpus. Note the order of the article: by q.~44 the phrase ``apart from God'' has been earned by qq.~2--11; the present essay reproduces that order in compressed form. The divine name appears in the texts quoted here; the essay's own identification of the terminus so named is not performed but deferred, and section~4 identifies where the tradition performs it.}

The same article's sequel closes the last door a reader might leave open. Perhaps existence is conferred on \emph{forms}, one might think, while the raw stuff underneath is simply given---the eternal, underived matter of the Greek picture. Aquinas denies it, and denies it precisely on the ground that a cause of being as being leaves nothing outside its effect: what causes things as beings must cause them ``not only according as they are `such' by accidental forms, nor according as they are `these' by substantial forms, but also according to all that belongs to their being at all in any way,'' and therefore ``it is necessary to say that also primary matter is created by the universal cause of things.''\endnote{\emph{Summa theologiae} I, q.~44, a.~2, corpus; checked 2026-07-25 as above. The reply to the third objection guards the claim's sense: the argument ``does not show that matter is not created, but that it is not created without form.'' Nothing in this essay depends on the hylomorphic analysis of bodies as such; the article is used for its principle, that a cause of being reaches whatever belongs to a thing's being in any way.} Nothing gets to stand outside. There is no floor beneath the floor.

At the head of the series of borrowers---and only there---stands something whose essence \emph{is} its existence: not one more thing that happens to exist, but subsistent existence itself, having nothing in it of the unfilled capacity that would need, in turn, to be filled. Aquinas's argument for the identification is threefold and each strand is short. Whatever a thing has besides its essence is caused either by its own principles or by an external agent; but nothing can cause its own existence; ``therefore that thing, whose existence differs from its essence, must have its existence caused by another''---which cannot be true of the first cause. Again, ``existence is that which makes every form or nature actual,'' so existence stands to essence as act to potency, and in the first there is no potency. And again, ``that which has existence but is not existence, is a being by participation''---and a participated being is not first.\endnote{\emph{Summa theologiae} I, q.~3, a.~4, corpus (\latin{Deus non solum est sua essentia\ldots{} sed etiam suum esse}); checked 2026-07-25, English at New Advent, Latin at Corpus Thomisticum. The three arguments summarized are the article's own, in its order. The second is the one that matters most for what follows: existence is ``that which makes every form or nature actual; for goodness and humanity are spoken of as actual, only because they are spoken of as existing.''} Aquinas's name for the terminus is blunt: in God alone essence and existence are identical; God is ``not only His own essence\ldots{} but also His own existence.''

\subsection{A name in the present tense}

One further passage deserves its place here, because it shows the tradition drawing from this analysis a consequence about language rather than about metaphysics. Asked which of the divine names is most proper, Aquinas answers: not ``God,'' and not ``the good,'' but ``He Who Is''---and he gives three reasons, of which the third is the one that concerns this essay. The name is most fitting ``from its consignification, for it signifies present existence; and this above all properly applies to God, whose existence knows not past or future.''\endnote{\emph{Summa theologiae} I, q.~13, a.~11, corpus; checked 2026-07-25 as above. The other two reasons given are that the name signifies existence itself rather than any determinate form, and that it is maximally undetermined---whence the article's quotation of John Damascene, that ``He Who Is'' ``contains existence itself as an infinite and indeterminate sea of substance.'' The article's reply to the first objection preserves the counterweight: as regards the object intended, the name ``God'' is more proper, and more proper still is the Tetragrammaton. This essay quotes the article for its tense, not to settle the theology of the divine names.} Present existence: the tense of the argument all the way down. The question of the first pages was never about a distant beginning, and the tradition's most careful sentence about the divine name says why the question could not have been about that. What subsistent existence has is not a long history of existing. It is existing, now, without the qualification ``still.''
